\begin{exercisechunk}
\pagebreak[3]
\item \textbf{EM increases likelihood but can stop at a saddle.}
\label[exercise]{ex:l5-em}
\exerciseprofile{3}{2}{3}{\exproof\quad\excalculation\quad\exconstruction\quad\exinterpretation}{%
Prove EM monotonicity and construct an exact EM iteration that remains at a
saddle point of the observed log likelihood.}
Let the unobserved variable \(X\) and observed variable \(Y\) take values
in finite sets. Fix an observed value \(y_{\mathrm{obs}}\), and assume that
\(Q_\theta(x,y_{\mathrm{obs}})>0\) for every allowed \(\theta,x\). Write
\(\ell(\theta)=\log Q_\theta(y_{\mathrm{obs}})\). Starting from \(\theta_k\), one EM
iteration consists of the E-step
\[
 q_k(x)=Q_{\theta_k}(X=x\mid Y=y_{\mathrm{obs}})
\]
and the M-step
\[
 \theta_{k+1}\in\arg\max_\theta
 \sum_x q_k(x)\log Q_\theta(x,y_{\mathrm{obs}}).
\]
A generalized M-step need only make this sum no smaller than its value at
\(\theta_k\). For part (a), fix \(\theta_0\) and abbreviate
\(q(x)=Q_{\theta_0}(X=x\mid Y=y_{\mathrm{obs}})\).
\begin{enumerate}[(a),beginpenalty=10000]
\item Prove
\[
 \ell(\theta)-\ell(\theta_0)
 =\sum_x q(x)\log\frac{Q_\theta(x,y_{\mathrm{obs}})}
                              {Q_{\theta_0}(x,y_{\mathrm{obs}})}
   +\KL(q,Q_\theta(\cdot\mid y_{\mathrm{obs}})).
\]
Deduce monotonicity of exact EM and of an M-step that merely increases
the expected complete-data log likelihood.
\item The proof in part (a) depends on \(q\) being the full current
posterior. Identify the equality that is lost if \(q\) is replaced by a
point mass at a most likely hidden value. Explain also why monotonicity of
the likelihood values does not guarantee convergence to a local maximum.
\item To exhibit the stationary-point problem from the second subproblem
of part (b), consider parameters \((a,b)\in(-1,1)^2\), with
\[
 Q_{a,b}(X=x)=\frac{1+ax}{2},\qquad
 Q_{a,b}(Y=1\mid X=x)=\frac{1+bx}{2},
 \qquad x\in\{-1,1\}.
\]
Suppose that \(Y=1\) is observed. Show that
\(\ell(a,b)=\log((1+ab)/2)\). Starting from \((a,b)=(0,0)\), compute
the E- and M-steps and show that exact EM remains at the origin. Use the
Hessian of \(\ell\) to show that the origin is a saddle point.
\item The stationary point in part (c) has an indefinite Hessian. Now
suppose that EM converges to an interior stationary point \(\theta_\star\).
Assume that the observed log likelihood is twice continuously
differentiable near \(\theta_\star\) and
\(\nabla^2\ell(\theta_\star)\) is negative definite. What can be
concluded about \(\theta_\star\)? Justify the conclusion.
\end{enumerate}
\begin{hints}
For (a), factor the joint mass as
\(Q_\theta(x,y_{\mathrm{obs}})
=Q_\theta(x\mid y_{\mathrm{obs}})Q_\theta(y_{\mathrm{obs}})\).
For (c), first compute the posterior of \(X\) at the origin. The
expected complete-data log likelihood then separates into a function of
\(a\) and a function of \(b\).
\end{hints}
\end{exercisechunk}
